\section{Installation}
\label{sec:install}

\subsection{Requirements}

\begin{itemize}[leftmargin=1.4em]
  \item \textbf{Python 3.13} or newer.  The project is pinned to 3.13 and the
        test suite asserts the interpreter version.
  \item \textbf{uv}, which manages the interpreter, the virtual environment and the
        development tools.
  \item No runtime dependency: the engine uses the standard library only.
\end{itemize}

\subsection{From a checkout}

\begin{lstlisting}[style=shell]
# fetch a managed Python 3.13 if you do not have one
uv python install 3.13

# create .venv, install the project and the development tools
uv sync

# the stage gate: formatting, linting and the whole test suite
scripts/check.sh
\end{lstlisting}

\texttt{uv sync} installs the project in editable mode and puts the
\texttt{rubiks-cfop} console script in \texttt{.venv/bin}.  You can run it directly,
through \texttt{uv run}, or through \texttt{uv tool install .} if you prefer a
standalone command.

\begin{lstlisting}[style=shell]
uv run rubiks-cfop --scramble "R U R' U' F2"
\end{lstlisting}

\subsection{Building a wheel}

\begin{lstlisting}[style=shell]
uv build            # dist/rubiks_cfop-0.1.0-py3-none-any.whl and .tar.gz
\end{lstlisting}

The wheel is self-contained: the JSON lookup tables live inside the package
(\texttt{rubiks\_cfop/data/}) and are installed with it, and the console script is
declared as an entry point.

\subsection{What is installed}

\begin{center}
\begin{tabular}{@{}ll@{}}
\toprule
Path & Contents \\
\midrule
\texttt{rubiks\_cfop/cube/} & facelets, moves, notation, \texttt{Cube}, \texttt{CubieCube} \\
\texttt{rubiks\_cfop/scramble/} & random-state and random-move generators \\
\texttt{rubiks\_cfop/solver/} & cross, F2L, OLL/PLL, CFOP pipeline, two-phase \\
\texttt{rubiks\_cfop/render/} & the unfolded net and its text serializers \\
\texttt{rubiks\_cfop/data/} & \texttt{oll.json}, \texttt{pll.json}, \texttt{f2l.json} \\
\texttt{rubiks\_cfop/cli.py} & the \texttt{rubiks-cfop} command \\
\bottomrule
\end{tabular}
\end{center}

\subsection{The first run is slower}

The random-state scrambler needs a set of precomputed tables: coordinates, move
tables and four pruning tables totalling about six megabytes.  They are built on
first use and then cached for the life of the process, which takes roughly ten
seconds of CPU.  Anything that runs the scrambler should build them up front on a
background thread:

\begin{lstlisting}[style=python]
import rubiks_cfop as cube

cube.preload()          # build the scrambler tables now
scramble = cube.generate_scramble(seed=1)
\end{lstlisting}

Solving a cube you already have does not need those tables, only the small cross
table (one megabyte, under a second).
